% ============ 前言(原书まえがき,原书 pp.3-8 = PDF pp.3-8) ============
\chapter*{前言}\addcontentsline{toc}{chapter}{前言}

调和分析一词在与正交级数展开相关的地方被广泛使用,但本书所处理的,是在某种意义下直接与三角级数或 Fourier 变换相关的古典调和分析.不过,所谓古典,是指题材是自古以来就被处理的,这一领域目前仍在被积极地研究,其成果和方法也被活跃地用于分析学的其他领域;此外,它还包含具有重要意义的一些未解决问题,是可能参与 21 世纪整个分析学的本质性进展的领域.

本书以在大学数学专业中掌握了复函数论,Lebesgue 积分和泛函分析的基本内容,并打算在研究生阶段专攻分析学的学生们为对象,其主要目的是引导年轻人进入古典调和分析的研究.同时,也期望已经从事分析学各领域研究的年轻人,能把这一领域的理论作为一种素养来熟悉.本书不是对古典调和分析的包罗万象的详细解说,而是作为对这一领域的邀请,解说若干话题.各章相互独立撰写.下面,对本书的内容作稍详细的说明.

第\ref{ch:1}章是对目前常被称为调和分析的实函数论方法的理论的解说.单变量的 $p$ 次可积函数的 Fourier 级数的部分和是否在 $L^{p}$ 范数意义下收敛于原函数的问题,可归结为被称为共轭函数变换或 Hilbert 变换的某个奇异积分算子的 $L^{p}$ 范数有界性.此外,对函数 $f$,把 Laplace 方程 $\Delta u = f$ 的解 $u$ 的二阶偏导数 $\partial^{2}u/\partial x_{i}\partial x_{j}$ 与之对应的线性算子,也是同类的奇异积分算子.这些奇异积分算子的 $L^{p}$ 理论,是 1950 年前后由 A.\,P.\ Calderón 和 A. Zygmund 开创的.作为这一理论的钥匙而被使用的,是分解函数的方法和把算子的弱型估计进行插值的方法,这些是实函数论的方法.实函数论的方法在 1970 年代由 C. Fefferman 和 E.\,M.\ Stein 等人发展成实函数论式的 $H^{p}$ 空间理论.这些成果和方法目前也被有效地用于偏微分方程论等领域.第\ref{ch:1}章就单变量函数解说这一实函数论方法.

第\ref{ch:2}章解说单变量全纯函数的 Hardy 空间理论.构成 Fourier 级数的函数 $e^{in\theta}$,在 $n \geq 0$ 时是全纯函数 $z^{n} = (re^{i\theta})^{n}$ 在 $r = 1$ 处的边界值.由此,把单位圆周上的函数与单位圆盘内的全纯函数对应起来,利用全纯函数的性质来研究 Fourier 级数的性质,这在 20 世纪前半叶曾被积极地研究,被称为 Fourier 级数的复解析方法.例如,单位圆周上的复值测度 $\mu$ 的 Fourier 系数 $\hat{\mu}(n)$ 对负的 $n$ 全为 0 时,$\mu$ 关于 Lebesgue 测度绝对连续----这个 Riesz 兄弟定理可用复解析方法简明地证明.此外,本书虽不涉及,但关于 Hilbert 变换的加权 $L^2$ 范数有界性的 Helson--Szeg\"o 定理,可利用单位圆周上满足 $\log w$ 可积的正值函数 $w$ 等于单位圆盘内全纯函数的边界值的绝对值这一事实,从复函数的性质推导出来.在这个复解析方法中扮演重要角色的全纯函数 Hardy 空间,作为一个函数空间的结构等,其本身提供了许多有趣的问题,得到了大量研究.关于 Hardy 空间各种形式推广的研究,即使在最近也仍是古典调和分析的重要研究课题.

第\ref{ch:3}章是 Fourier 分析中 Banach 环理论的解说.这一理论的发端是 N. Wiener 定理:具有绝对收敛 Fourier 级数的连续函数,若处处不为零,则其倒函数仍具有绝对收敛的 Fourier 级数.这一结果在 Fourier 分析中提供了两个话题.其一是谱综合问题,例如,对于具有绝对收敛 Fourier 级数的连续函数的空间 $A(\mathbf{T})$,可以表述为:"对单位圆的任意闭集,在其上为 0 的 $A(\mathbf{T})$ 的函数,能否由在该闭集邻域内为 0 的 $A(\mathbf{T})$ 的函数在范数意义下逼近(谱综合集)?"此后,这一问题经 I. Gelfand 的理论被推广到交换 Banach 环.1948 年,L. Schwartz 证明了 3 维以上 Euclid 空间的单位球面不是 $A(\mathbf{R}^{n})$ 的谱综合集.决定性的结果是 1959 年 P. Malliavin 关于"$A(\mathbf{T})$ 中非谱综合集的存在性".在与 Fourier 变换有关的 Banach 环中,他们的研究是结合几何性质以及(Lebesgue 测度为零的)稀疏集的性质发展起来的.其二是被称为作用函数的问题,例如,对 $A(\mathbf{T})$:"$[-1,1]$ 上的怎样的函数作用于 $A(\mathbf{T})$?"1930 年代初期,N. Wiener 与 P. Levy 证明了"解析函数作用于 $A(\mathbf{T})$",而 1958 年 Y. Katznelson 证明了其逆也成立.在 P. Malliavin 与 Y. Katznelson 的两个定理中,证明的本质部分从互为共轭的立场进行范数估计,这一点非常有趣.作用函数问题也被推广到交换 Banach 环并得到发展.谱综合与作用函数的研究,是通过考察交换 Banach 环的一个特征而发展起来的;而最近,本书虽不涉及,但已结合调和分析,在各种函数空间中开展研究.

第\ref{ch:4}章是对被称为 Kakeya 问题的,古典 Fourier 分析中若干重要未解决问题的解说.Kakeya 问题,是关于多变量函数或高维图形的一些互相关联的问题的总称.例如,确定多变量 Fourier 级数的被称为 Bochner--Riesz 平均的求和法在 $L^{p}$ 范数意义下收敛的 $p$ 的范围的问题;又如,对高维 Euclid 空间 $\mathbb{R}^{d}$ 上的函数 $f$,把 $f$ 的 Fourier 变换 $\hat{f}$ 到球面 $S$ 的限制 $\hat{f}|S$ 与之对应的算子 $f \mapsto \hat{f}|S$ 成为从 $L^{p}(\mathbb{R}^{d})$ 到 $L^{q}(S)$ 的有界算子的指数 $(p,q)$ 的确定问题.近年来,已明确这些问题与被称为 Nikodym 极大算子的 $\mathbb{R}^{d}$ 上算子的 $L^{p}$ 有界性问题,被称为 Kakeya 集的 $\mathbb{R}^{d}$ 子集的 Hausdorff 维数估计问题等深刻相关.这些问题在二维情形均已解决,但在三维以上情形,与其表面上的简单相反,已弄清楚是非常困难的问题.近年 Wolff 和 Bourgain 等人取得了重要成果,但尚未达到解决.这些问题的核心常常归结为估计满足某条件的有限集合的元素个数的问题,对此的讨论被称为组合论证.这些问题得到解决,并且其中组合论证作为可理解的东西被纳入分析学,尚有待将来.

本书的内容就是以上 4 章,但作为对古典调和分析的邀请,也许还可以设立一章关于概率论的.关于间隙 Fourier 级数与独立随机变量序列的相似性,自古以来就有大量研究.Burkholder,Gundy 与 Silverstein 首次给出用 Brown 运动定义的最大函数刻画全纯函数 Hardy 空间中函数的实部的定理时,该最大函数成为关键.另外,内山明人用实函数论方法证明 BMO 的 Fefferman--Stein 分解,其最初是在鞅的设定中进行的.奇异积分,实函数论 Hardy 空间,BMO 空间等,与对实变函数并行的理论一样,也可在概率论的鞅形式下建立.关于这些,请参考以下文献等.

[1] D. L. Burkholder, R. F. Gundy and M. L. Silverstein, A maximal function characterization of the class $H^{p}$, Trans. Amer. Math. Soc. 157 (1971), 137--153.

[2] A. M. Garsia, Martingale Inequalities, Seminar Notes on Recent Progress, Mathematics Lecture Note Series, W. A. Benjamin, 1973.

[3] K. E. Petersen, Brownian Motion, Hardy Spaces and Bounded Mean Oscillation, London Mathematical Society Lecture Note Series 28, Cambridge University Press, 1977.

[4] D. L. Burkholder, Martingale theory and harmonic analysis in Euclidean spaces, in Harmonic Analysis in Euclidean Spaces, G. Weiss and S. Wainger eds., part 2, Proc. Symp. Pure Math. Vol. 35, American Math. Soc., 1979, pp. 283--301.

[5] A. Uchiyama, A constructive proof of the Fefferman--Stein decomposition of BMO on simple martingales, Conference on Harmonic Analysis in honor of Antoni Zygmund, Volume II, W. Beckner et al. eds, pp. 495--505, 1983.

各章的执笔:第\ref{ch:1}章为薮田公三,第\ref{ch:2}章为中路貴彦,第\ref{ch:3}章为佐藤圓治,第\ref{ch:4}章为田中仁.全书由宮地晶彦负责统稿.

中井英一先生与富田直人先生通读了第\ref{ch:1}章,井上純治先生与山本隆範先生通读了第\ref{ch:2}章,均从原稿阶段起仔细阅读并提出了各种宝贵意见.另外,朝倉書店编辑部从本书的策划阶段起给予了多方面的帮助.在此谨表深切谢意.

2008 年 2 月

宮地晶彦
